\documentclass[12pt]{article}
\usepackage{amsmath,amsthm,amssymb,mathtools}
\usepackage{hyperref}

\newtheorem{theorem}{Theorem}[section]
\newtheorem{lemma}[theorem]{Lemma}
\newtheorem{proposition}[theorem]{Proposition}
\newtheorem{corollary}[theorem]{Corollary}
\newtheorem{definition}[theorem]{Definition}
\newtheorem{remark}[theorem]{Remark}
\newtheorem{obligation}{Obligation}[section]

\newcommand{\R}{\mathbb{R}}
\newcommand{\Z}{\mathbb{Z}}
\DeclareMathOperator{\Sym}{Sym}

\title{An Explicit Constant from the GL(3) Voronoi Formula\\
\large Proof Document 05 --- sym2-effective-lbound\\
\normalsize \textbf{Status: [OBL]}}
\author{}
\date{2026-08-18}

\begin{document}
\maketitle

\begin{abstract}
This document records the precise obligation needed to turn the Miller--Schmid
GL(3) Voronoi summation formula into an explicit constant for the dual/contour
term in the AFE approach to \(L(1,\mathrm{sym}^2\Delta)\), and eventually into
the zero-free-region estimate M-3.  No explicit constant is proved here.  The
only admitted input is the exact normalized formula in Miller--Schmid Theorem
1.18; numerical thresholds obtained by floating-point scans are discovery-tier
evidence and have no logical force.
\end{abstract}

\section{Source formula}

\begin{theorem}[Miller--Schmid, Theorem 1.18 {\normalfont [BASE]}]
\label{thm:ms-voronoi}
Let \(a_{n,m}\) be Fourier coefficients of a cuspidal \(GL(3,\Z)\)-automorphic
representation with parameters \(\lambda=(\lambda_1,\lambda_2,\lambda_3)\) and
\(\delta=(\delta_1,\delta_2,\delta_3)\).  Let \(f\) be a Schwartz function which
vanishes to infinite order at the origin.  For \((a,c)=1\), \(c\neq0\), an inverse
\(\bar a\) of \(a\bmod c\), and \(q>0\),
\begin{align*}
  \sum_{n\neq0} a_{q,n} e(-na/c) f(n)
  &= \sum_{d\mid cq} \left|\frac{c}{d}\right|
       \sum_{n\neq0} \frac{A(n,d)}{|n|}
       S(q\bar a,n;qc/d)\,
       F\!\left(\frac{nd^2}{c^3q}\right).
\end{align*}
Here \(S\) is the Kloosterman sum and \(F\) is the normalized threefold
integral transform specified in the source.  In particular, the factor
\(A(n,d)/|n|\), the modulus \(qc/d\), and the argument
\(nd^2/(c^3q)\) are part of the theorem's normalization.
\end{theorem}

\begin{remark}
The source-backed entry `MS-V.1` in
\texttt{baseline/REFERENCE\_BASELINE.md} records the public primary source.
This repository has not independently reproved Theorem
\ref{thm:ms-voronoi}; it is admitted only as the stated baseline.  Any derivation
which replaces the displayed normalization by a generic expression such as
\(c^{-2}K_\nu(\cdot)\) is not a specialization of this theorem.
\end{remark}

\section{Required explicit estimate}

Let \(\mathcal{T}\) denote a fixed finite-dimensional test-function class to be
specified in the eventual certificate.  It must include the Mellin--Bessel weight
used by the \(J\)-term, together with explicit Sobolev or Schwartz seminorms.

\begin{obligation}[Explicit GL(3) Voronoi constant {\normalfont [OBL]}]
\label{ob:explicit-voronoi}
Derive, from Theorem~\ref{thm:ms-voronoi}, an explicit function
\[
  C_{\mathrm{Vor}}\colon \mathcal{T}\to(0,\infty)
\]
such that for every \(X\ge X_0\), every \(q\) and \(c\) in the finite set
occurring in the \(J\)-term, and every \(\phi\in\mathcal{T}\),
\[
  \left|
    \sum_{n\neq0} a_{q,n} e(-na/c)\phi(n/X)
  \right|
  \le C_{\mathrm{Vor}}(\phi) X^{2/3+\epsilon}
\]
(or prove a different explicit exponent rigorously).  All constants, including
\(X_0\), \(\epsilon\), the coefficient moment, and the dependence on \(q,c\),
must be computable and replayable.  No numerical threshold from
\texttt{discovery/} may be substituted for \(C_{\mathrm{Vor}}\).
\end{obligation}

\begin{remark}
For the self-dual symmetric-square lift, replacing a primal sum by the same sum
with \(q=c=1\) is tautological and cannot prove cancellation.  The derivation
must explain why the chosen additive characters and dyadic decomposition convert
the left side into shorter oscillatory dual sums.  A mere identity with a dual
sum is not a size estimate.
\end{remark}

\subsection{Required constant-tracking ledger}

Obligation~\ref{ob:explicit-voronoi} is closed only after all following entries
are supplied and independently checked:

\begin{enumerate}
  \item \textbf{Test function normalization.} An explicit definition of
        \(\mathcal{T}\), its Mellin transform, and every seminorm used in the
        bound.  The \(J\)-weight may not be changed after the estimate is proved.
  \item \textbf{Coefficient normalization.} A definition of \(a_{q,n}\) and
        \(A(n,d)\) matching Miller--Schmid.  The transition from normalized
        \(\Delta\) coefficients to \(a_{q,n}\) must be an identity, not a
        heuristic.
  \item \textbf{Rankin--Selberg moment.} A source-backed explicit bound for the
        required finite \(A(n,d)\)-moment, including the diagonal range and all
        powers of \(d\).
  \item \textbf{Kloosterman estimate.} The chosen Weil-type bound, its exact
        constant, and its dependence on \(qc/d\).
  \item \textbf{Bessel transform norm.} A certified upper bound for the relevant
        \(L^1\), \(L^2\), or Sobolev norm of \(F\), using outward-rounded
        interval arithmetic.
  \item \textbf{Summation order.} A rigorous dyadic decomposition of the
        \((c,d,n)\)-sum.  The certificate must state the resulting \(c\)-sum
        exponent; \(\zeta(3/2)\), \(\zeta(4/3)\), and \(\zeta(7/6)\) are distinct
        candidate calculations, not interchangeable.
  \item \textbf{Truncation and tails.} Explicit truncation radii and tail bounds
        for every infinite or conditionally convergent sum.  If cancellation is
        used, it must be bounded rather than observed numerically.
  \item \textbf{Replay command.} A deterministic checker command which recomputes
        every numerical constant from exact or interval inputs.
\end{enumerate}

\section{Forbidden shortcuts}

\begin{itemize}
  \item An empirical maximum of finitely many sums is not \(C_{\mathrm{Vor}}\).
  \item A numerically large margin below a floating-point threshold does not
        close M-3 or the \(J\)-certification.
  \item The theorem is an exact summation formula, not by itself an estimate;
        cancellation must come from the Kloosterman, coefficient, and Bessel
        inputs.
  \item No bound may silently truncate the \(d\mid cq\), \(c\), or \(n\)-sums.
  \item The output may not be called a zero-free region until it is connected
        to a complete certified argument for \(L(s,\mathrm{sym}^2\Delta)\).
\end{itemize}

\begin{thebibliography}{9}
\bibitem{MS06} S.~D.~Miller and W.~Schmid,
  \emph{Automorphic distributions, \(L\)-functions, and Voronoi summation for
  \(GL(3)\)}, Ann.\ of Math.\ (2) \textbf{164} (2006), no.~2, 423--488.
\end{thebibliography}

\end{document}
